\documentclass[a4paper]{article}
\usepackage[francais]{babel}
\usepackage[T1]{fontenc}  
\usepackage[utf8x]{inputenc} 
\usepackage{amsmath}
\usepackage{amsfonts}
\usepackage{amssymb}
\usepackage{graphicx}
\usepackage[a4paper, left=2cm, right=2cm, top=2.5cm, bottom=2cm, headheight=40pt]{geometry}
\usepackage{color}
\usepackage{mathtools}
\usepackage{algcompatible}
\usepackage{algorithm}
\usepackage{algorithmicx}
\usepackage{algpseudocode}
\usepackage[dvipsnames]{xcolor}
\usepackage{tikz}
\usetikzlibrary{shapes}
\usepackage[square,sort,comma,numbers]{natbib}
%\usepackage{emptypage}
\usepackage{hyperref}
\hypersetup{
    colorlinks,
    citecolor=black,
    filecolor=black,
    linkcolor=black,
    urlcolor=black
}
\newcount\colveccount
\newcommand*\colvec[1]{
        \global\colveccount#1
        \begin{pmatrix}
        \colvecnext
}
\def\colvecnext#1{
        #1
        \global\advance\colveccount-1
        \ifnum\colveccount>0
                \\
                \expandafter\colvecnext
        \else
                \end{pmatrix}
        \fi
}
\title{Hashing Functions}
\author{}


\begin{document}
\begin{titlepage}
  \begin{sffamily}
  \begin{center}

    Univesité de Bordeaux\\ Sciences \& Technologie\\
           351 Cours de la Liberation\\33400 Talence\\[1em]
            \textbf{\underline{- Projet Cryptanalyse -}}\\[1.5cm]
           
    \includegraphics[scale=0.11]{UB.jpg}
    \\[3cm]

    { \huge \bfseries Projet sur la cryptanalyse de A5/2\\[0.5cm] }
    \begin{flushright}
      \bfseries {Amelie GUEMON \& Ida TUCKER\\[1em]Master 2 CSI --- Cryptologie \& Sécurité Informatique}\\[6.6cm]
    \end{flushright}
    
   
    \today

  \end{center}
  \end{sffamily}
\end{titlepage}

%% QUESTION 1
\section{Implémentation d du chiffrement A5/2 }
Implémentation du chiffrement A5/2, cf. a5-2.sage

%% QUESTION 2
\section{Expression des registres au moyen d’équations linéaires: théorique}


Avant le $1^{er}$ passage par la fonction A5/2-step on a: 
\begin{itemize}
\item $R_1 = (x_0, \ldots, x_{18})$
\item $R_2 = (x_{19}, \ldots, x_{40})$
\item $R_1 = (x_{41}, \ldots, x_{63})$
\end{itemize}

Notons, au $k^{ième}$ passage par la fonction A5/2-step:
\begin{itemize}
\item $R_1 = (x^{\lbrack k\rbrack}_0, \ldots, x^{\lbrack k \rbrack}_{18})$
\item $R_2 = (x^{\lbrack k \rbrack}_{19}, \ldots, x^{\lbrack k \rbrack}_{40})$
\item $R_3 = (x^{\lbrack k \rbrack}_{41}, \ldots, x^{\lbrack k \rbrack}_{63})$
\end{itemize}

Montrons qu'on peut exprimer le contenu des registres $R_1$, $R_2$, $R_3$ au moyen d'équations linéaires en les $x_i$ par récurrence.

Il est clair qu'à l'étape zero (avant le premier passage par A5/2-step), la propriété est vraie.
Supposons la vraie au rang $k$ pour $k \geq 0$.
On a donc $x^{\lbrack k\rbrack}_0, \ldots, x^{\lbrack k \rbrack}_{63}$ qui s'expriment comme équation linéaire des $x_i$.
On passe alors dans la fonction A5/2-step:
\begin{itemize}
\item Si m = $R_{4,6}$ on met à jour le $LFSR_1$, et $R_1$ devient:
$$R_1 = (x^{\lbrack k+1\rbrack}_0, \ldots, x^{\lbrack k+1 \rbrack}_{18}) = (x^{\lbrack k\rbrack}_1, x^{\lbrack k\rbrack}_2, \ldots, x^{\lbrack k \rbrack}_{18}, x^{\lbrack k\rbrack}_0 + x^{\lbrack k\rbrack}_1 + x^{\lbrack k\rbrack}_2 + x^{\lbrack k\rbrack}_5)$$
\item Si m = $R_{4,13}$ on met à jour le $LFSR_2$, et $R_2$ devient:
$$R_2 = (x^{\lbrack k+1\rbrack}_{19}, \ldots, x^{\lbrack k+1 \rbrack}_{40}) = (x^{\lbrack k\rbrack}_{20},x^{\lbrack k\rbrack}_{21}, \ldots, x^{\lbrack k+1 \rbrack}_{40}, x^{\lbrack k\rbrack}_{19} + x^{\lbrack k\rbrack}_{20})$$
\item Si m = $R_{4,9}$ on met à jour le $LFSR_3$, et $R_3$ devient:
$$R_3 = (x^{\lbrack k+1\rbrack}_{41}, \ldots, x^{\lbrack k+1 \rbrack}_{63}) = (x^{\lbrack k\rbrack}_{42}, x^{\lbrack k\rbrack}_{43}, \ldots, x^{\lbrack k\rbrack}_{63},x^{\lbrack k\rbrack}_{41} + x^{\lbrack k\rbrack}_{42}+ x^{\lbrack k\rbrack}_{43}+ x^{\lbrack k\rbrack}_{56})$$
\end{itemize}

On a donc le dernier élément de chaque registre au $(k+1)^{ième}$ passage qui est une somme des contenus des registre au $k^{ième}$ passage.
C'est bien une équation linéaire en les $x_i$ par linéarité de la somme.
Les autres éléments des registres sont les éléments de l'étape précédente décalés. C'est à dire $x^{\lbrack k+1\rbrack}_{i} = x^{\lbrack k\rbrack}_{i+1}$ pour $ i \neq 18, 40, 63$.

On a bien montré que durant toutes les étapes de la production de suite chiffrante de A5/2, on peut exprimer les contenus des registres $R_1$, $R_2$, $R_3$ au moyen d'équations linéaires en les $x_i$. 


%% QUESTION 1
\section{Expression des registres au moyen d'équations linéaires}
Expression des Registres $R_1$, $R_2$, $R_3$ au moyen d'équations linéaires en les $x_i$ en sage, cf. a5-2.sage

%% QUESTION 4
\section{Expression des bits de z comme équations quadratiques en les $x_i$: théorique}
\begin{table}[h!]
\centering
\begin{tabular}{l l l|c}
a & b & c & Maj(a,b,c) \\\hline
0 & 0 & 0 & 0  \\
0 & 0 & 1 & 0 \\
0 & 1 & 0 & 0 \\
0 & 1 & 1 & 1 \\
1 & 0 & 0 & 0  \\
1 & 0 & 1 & 1 \\
1 & 1 & 0 & 1 \\
1 & 1 & 1 & 1 \\

\end{tabular}
\caption{\label{tab:widgets}Table de vérité de Maj}
\end{table}

On voit, d'après sa table de vérité, que la forme algébrique normale de la fonction booléenne Maj est:
$$f(a,b,c) = ab \oplus ac \oplus bc$$
\paragraph{}
Au $(99+i)^{ième}$ passage par la fonction A5/2-step dans la production, on produit le $(i)^{ième}$ bit de la suite chiffrante: 

\begin{equation}
\begin{aligned}
  &z_i =& x^{\lbrack i + 99\rbrack}_0 + \mbox{Maj}(x^{\lbrack i+99\rbrack}_{3}, x^{\lbrack i+99\rbrack}_{4} + 1, x^{\lbrack i+99\rbrack}_{6})) \\
  & &+ x^{\lbrack i + 99\rbrack}_{19} + \mbox{Maj}(x^{\lbrack i+99\rbrack}_{27}, x^{\lbrack i+99\rbrack}_{24} + 1, x^{\lbrack i+99\rbrack}_{31})) \\
  & &+ x^{\lbrack i + 99\rbrack}_{41} + \mbox{Maj}(x^{\lbrack i+99\rbrack}_{45}, x^{\lbrack i+99\rbrack}_{50} + 1, x^{\lbrack i+99\rbrack}_{47})) \\
\end{aligned}
\end{equation}

On peut remplaçer la fonction Maj par sa forme normale algébrique et on obtient:

\begin{equation}
\begin{aligned}
  &z_i =& x^{\lbrack i + 99\rbrack}_0 + (x^{\lbrack i+99\rbrack}_{3}( x^{\lbrack i+99\rbrack}_{4} + 1) + x^{\lbrack i+99\rbrack}_{3} x^{\lbrack i+99\rbrack}_{6} + ( x^{\lbrack i+99\rbrack}_{4} + 1) x^{\lbrack i+99\rbrack}_{6}) \\  
  & &+ x^{\lbrack i + 99\rbrack}_{19} + (x^{\lbrack i+99\rbrack}_{27}(x^{\lbrack i+99\rbrack}_{24} + 1) + x^{\lbrack i+99\rbrack}_{27} x^{\lbrack i+99\rbrack}_{31} + (x^{\lbrack i+99\rbrack}_{24} + 1)x^{\lbrack i+99\rbrack}_{31}) \\  
  & &+ x^{\lbrack i + 99\rbrack}_{41} + (x^{\lbrack i+99\rbrack}_{45} (x^{\lbrack i+99\rbrack}_{50} + 1) + x^{\lbrack i+99\rbrack}_{45} x^{\lbrack i+99\rbrack}_{47} + (x^{\lbrack i+99\rbrack}_{50} + 1)x^{\lbrack i+99\rbrack}_{47}) \\
\end{aligned}
\end{equation}


On a vu que, quelle que soit l'étape k considérée dans la production, $ \forall i \in \left\{ {0,1 \ldots 63}\right\}$ les $x^{\lbrack k \rbrack}_{i}$ sont des équations linéaires des $x_i$. De plus dans le calcul des bits de suite chiffrante détaillé ci-dessus, on a uniquement des additions et des multiplications d'ordre 2 des $x^{\lbrack k \rbrack}_{i}$. Ceci montre bien que les bits de $z$ peuvent s'exprimer par des équations quadratiques en les $x_i$.


%% QUESTION 5
\section{Expression des équations quadratiques}
Expression des équations quadratiques en les $x_i$ des bits de la suite chiffrante en sage, cf. a5-2.sage

%% QUESTION 6
\section{Question 6}


%% QUESTION 7
\section{Linéarisation des équations}
Linéarisation des équations en sage, cf. a5-2.sage

%% QUESTION 8
\section{Application}
Application de la linéarisation des équations en sage sur une suite de 700 bits, cf. a5-2.sage

%% QUESTION 9
\section{Expression des registres en fonction de la clef et de l'IV: théorique}


Soit $A_1$ la matrice de rétroaction du $LFSR_1$. 
Montrons par récurrence que lors de l'étape d'initialisation, après la première boucle for, l'état du $LFSR_1$ est égal à
\begin{equation*}
X = \sum_{i=0}^{63} A_1^i \colvec{4}{0}{\vdots}{0}{K_{63-i}} \\
  = \colvec{4}{0}{\vdots}{0}{K_{63}} + A_1 \colvec{4}{0}{\vdots}{0}{K_{62}} + \cdots + A_1^{63}  \colvec{4}{0}{\vdots}{0}{K_{0}}
\label{eq:sn}
\end{equation*}

Après le premier passage par la première boucle for on a:

\begin{equation*}
X = \colvec{4}{0}{\vdots}{0}{K_{0}} \\
\label{eq:sn}
\end{equation*}

Après le deuxième passage par la première boucle for on a:
\begin{equation*}
X = A_1 \colvec{4}{0}{\vdots}{0}{K_{0}} + \colvec{4}{0}{\vdots}{0}{K_{1}} \\
= \sum_{i=0}^{1} A_1^i \colvec{4}{0}{\vdots}{0}{K_{1-i}}
\label{eq:sn}
\end{equation*}

Supposons pour $k \geq 1$ qu'au $k^{ieme}$ passage par la boucle on ai:
\begin{equation*}
X = \sum_{i=0}^{k} A_1^i \colvec{4}{0}{\vdots}{0}{K_{k-i}}\\
\label{eq:sn}
\end{equation*}

Alors au passage suivant par la boucle on a:

\begin{align*}
X &= A_1 \sum_{i=0}^{k} A_1^i \colvec{4}{0}{\vdots}{0}{K_{k-i}} + \colvec{4}{0}{\vdots}{0}{K_{k+1}} = \sum_{i=1}^{k+1} A_1^i \colvec{4}{0}{\vdots}{0}{K_{k+1-i}} + \colvec{4}{0}{\vdots}{0}{K_{k+1}} \\
X &= \sum_{i=0}^{k+1} A_1^i \colvec{4}{0}{\vdots}{0}{K_{k+1-i}}
\label{eq:sn}
\end{align*}

On obtient donc bien par récurrence que 
\begin{equation*}
X = \sum_{i=0}^{63} A_1^i \colvec{4}{0}{\vdots}{0}{K_{63-i}} \\
\label{eq:sn}
\end{equation*}


De même, le premier passage par la deuxième boucle for on a:
\begin{equation*}
R_1 = A_1 X + \colvec{4}{0}{\vdots}{0}{IV_{0}} = A_1^1 X + \sum_{i=0}^{0} A_1^i \colvec{4}{0}{\vdots}{0}{IV_{0-i}} \\
\label{eq:sn}
\end{equation*}

Pour $k \geq 0 $, on suppose:

\begin{equation*}
R_1 = A_1^{k+1} X + \sum_{i=0}^{k} A_1^i \colvec{4}{0}{\vdots}{0}{IV_{k-i}} \\
\label{eq:sn}
\end{equation*}

A l'étape $k+1$ on a alors:
\begin{align*}
R_1 &= A_1 \times ( A_1^{k+1} X + \sum_{i=0}^{k} A_1^i \colvec{4}{0}{\vdots}{0}{IV_{k-i}}) + \colvec{4}{0}{\vdots}{0}{IV_{k+1}} \\
&= A_1^{k+2} X + \sum_{i=0}^{k} A_1^{i+1} \colvec{4}{0}{\vdots}{0}{IV_{k-i}} + \colvec{4}{0}{\vdots}{0}{IV_{k+1}} \\
&= A_1^{k+2} X + \sum_{i=0}^{k+1} A_1^{i} \colvec{4}{0}{\vdots}{0}{IV_{k+1-i}}
\label{eq:sn}
\end{align*}

On obtient bien par récurrence qu'après la deuxième boucle for, l'état du $LFSR_1$ est

\begin{equation*}
A_1^{22} X + \sum_{i=0}^{21} A_1^i \colvec{4}{0}{\vdots}{0}{IV_{21-i}} \\
\label{eq:sn}
\end{equation*}


%% QUESTION 10
\section{Expression des registres en fonction de la clef et de l'IV: pratique}

En notant $A_2$, $A_3$ et $A_4$ les matrices de rétroactions des $LFSR_2$, $LFSR_3$ et $LFSR_4$, on trouve de même qu'après la première boucle for, les états respectifs des $LFSR_2$, $LFSR_3$ et $LFSR_4$ sont 
\begin{align*}
X_2 &= \sum_{i=0}^{63} A_2^i \colvec{4}{0}{\vdots}{0}{K_{63-i}} \\
X_3 &= \sum_{i=0}^{63} A_3^i \colvec{4}{0}{\vdots}{0}{K_{63-i}} \\
X_4 &= \sum_{i=0}^{63} A_4^i \colvec{4}{0}{\vdots}{0}{K_{63-i}} \\
\label{eq:sn}
\end{align*}

Après la deuxième boucle for, leurs états respectifs sont:

\begin{align*}
A_2^{22} X_2 + \sum_{i=0}^{21} A_2^i \colvec{4}{0}{\vdots}{0}{IV_{21-i}} \\
A_3^{22} X_3 + \sum_{i=0}^{21} A_3^i \colvec{4}{0}{\vdots}{0}{IV_{21-i}} \\
A_4^{22} X_4 + \sum_{i=0}^{21} A_4^i \colvec{4}{0}{\vdots}{0}{IV_{21-i}}
\end{align*}


%% QUESTION 11
\section{Expression des registres produisant $z_1$ et $z_2$ en fonction des $x_i$}
Notons $R_{i,j}$ l'état du registre i après la phase d'initialisation ayant produit les suites $z_j$.\\

Puisque l'IV de $z_0$ est nul, et d'après les questions 9 et 10 on a
\begin{align*}
R_{1,0} &= A_1^{22} X \\
R_{2,0} &= A_2^{22} X_2\\
R_{3,0} &= A_3^{22} X_3\\
R_{4,0} &= A_4^{22} X_4
\end{align*}
Pour $z_1$ seul $IV_{21}$ est non null, donc
\begin{align*}
R_{1,1} = A_1^{22} X + \colvec{4}{0}{\vdots}{0}{1} = R_{1,0} + \colvec{4}{0}{\vdots}{0}{1}\\
R_{2,1} = A_2^{22} X_2 + \colvec{4}{0}{\vdots}{0}{1} = R_{2,0} + \colvec{4}{0}{\vdots}{0}{1} \\
R_{3,1} = A_3^{22} X_3 + \colvec{4}{0}{\vdots}{0}{1} = R_{3,0} + \colvec{4}{0}{\vdots}{0}{1}\\
R_{4,1} = A_4^{22} X_4 + \colvec{4}{0}{\vdots}{0}{1} = R_{4,0} + \colvec{4}{0}{\vdots}{0}{1}
\end{align*}
D'où
\begin{align*}
R_{1,1} &= (x_0, \cdots ,x_{17}, x_{18}+1)\\
R_{2,1} &= (x_{19}, \cdots ,x_{39}, x_{40}+1)\\
R_{3,1} &= (x_{41}, \cdots ,x_{17}, x_{63}+1)\\
R_{4,1} &= R_{4,0} + \colvec{4}{0}{\vdots}{0}{1}
\end{align*}
Et enfin 
\begin{align*}
R_{1,2} &= R_{1,0} + A_1 \colvec{5}{0}{\vdots}{0}{0}{1} 
        = R_{1,0} + \colvec{6}{0}{\vdots}{0}{0}{1}{0} 
        = (x_0, \cdots ,x_{16},x_{17}+1, x_{18})\\
R_{2,2} &= R_{2,0} + A_2\colvec{5}{0}{\vdots}{0}{0}{1} 
        = R_{2,0} + \colvec{6}{0}{\vdots}{0}{0}{1}{0}
        = (x_{19}, \cdots ,x_{38},x_{37}+1, x_{40})\\
R_{3,2} &= R_{3,0} + A_3 \colvec{5}{0}{\vdots}{0}{0}{1} 
        = R_{3,0} + \colvec{6}{0}{\vdots}{0}{0}{1}{0}
        = (x_{41}, \cdots ,x_{61},x_{62}+1, x_{63}) \\
R_{4,2} &= R_{4,0} + A_4 \colvec{5}{0}{\vdots}{0}{0}{1}= 
          R_{4,0} + \colvec{6}{0}{\vdots}{0}{0}{1}{0}
\end{align*}



%% QUESTION 12
\section{Attaque permettant de retrouver la clef K à partir de $z_0, z_1, z_2$}
Afin de trouver les $x_i$ on aura, comme dans la question 8, à créer trois matrices (une pour chacune des suites $z_i$)
chacune possédant 228 équations linéaires en les 655 monomes (on suppose toujours $R_4$ connu, ainsi que les bits fixés à la fin de l'initialisation).
On a donc $684 = 3 \times 228$ équations linéaires et 655 inconnues, on peut résoudre le système et retrouver les $(x_i)_{0\leq i \leq 63}$.

Connaissant ces $x_i$ on pourra retrouver la clef, comme dans la question 9, en résolvant un système linéaire de $81$ équations à 64 inconnues.

En utilisant seulement l'IV produisant la suite $z_0$ on a 19 équations données par le registre $R_1$,
22 équations données par le registre $R_2$, 23 équations données par le registre $R_3$ et 17 équations données par le registre $R_4$.
Donc au total 81 équations linéairement dépendantes des bits de la clef K.

La résolution de ce système nous permet de retrouver la clef.


\nocite{*}


\end{document}
